\section{Limitations}
\label{sec:limitations}

The Shopify study has forty title-normalized decks and sixteen held-out product
clusters, but includes only the recovered complete prefix of thirty-nine of forty-eight
registered deep queries. The matching rule was changed after collection because the
original rule returned no decks. It is therefore exploratory. Seller domains do not
establish independent companies, and title normalization does not establish exact
variants or comparable shipping, taxes, returns, inventory, or checkout prices.

The recalled model assumes that a retained offer remains actionable at its observed
price. Neither source verifies this. UCP panels observe catalogs, not checkout, and
their lifecycle analysis is deliberately limited to observability quality. Shopify
Global Catalog coverage is not universal; an aggregator's visible graph is not the
set of all sellers a shopper could reach.

Inspection costs are a declared basis-point sensitivity grid, not measured billing or
user time. Seller prices within a product can be correlated, and random replay order
is a counterfactual ordering device rather than a record of user traffic. These are
reasons to preregister a larger collection, not reasons to reinterpret an inconclusive
interval as evidence of equivalence.